\section{Affine and quasi-affine schemes}
\label{section-affine-quasi-affine}

\begin{lemma}
\label{lemma-quasi-affine-O-ample}
Let $X$ be a scheme.
Then $X$ is quasi-affine if and only if $\mathcal{O}_X$ is ample.
\end{lemma}

\begin{proof}
Suppose that $X$ is quasi-affine. Set $A = \Gamma(X, \mathcal{O}_X)$.
Consider the open immersion
$$
j : X \longrightarrow \Spec(A)
$$
from Lemma \ref{lemma-quasi-affine}. Note that
$\Spec(A) = \text{Proj}(A[T])$, see
Constructions, Example \ref{constructions-example-trivial-proj}.
Hence we can apply Lemma \ref{lemma-open-in-proj-ample}
to deduce that $\mathcal{O}_X$ is ample.

\medskip\noindent
Suppose that $\mathcal{O}_X$ is ample.
Note that $\Gamma_*(X, \mathcal{O}_X) \cong A[T]$
as graded rings. Hence the result follows from Lemmas
\ref{lemma-ample-immersion-into-proj} and \ref{lemma-quasi-affine}
taking into account that
$\Spec(A) = \text{Proj}(A[T])$ for any ring $A$
as seen above.
\end{proof}

\begin{lemma}
\label{lemma-quasi-affine-locally-closed}
Let $X$ be a quasi-affine scheme. For any quasi-compact immersion
$i : X' \to X$ the scheme $X'$ is quasi-affine.
\end{lemma}

\begin{proof}
This can be proved directly without making use of the material on
ample invertible sheaves; we urge the reader to do this on a napkin.
Since $X$ is quasi-affine, we have that $\mathcal{O}_X$ is ample by
Lemma \ref{lemma-quasi-affine-O-ample}.
Then $\mathcal{O}_{X'}$ is ample by
Lemma \ref{lemma-ample-on-locally-closed}. Then $X'$ is quasi-affine by
Lemma \ref{lemma-quasi-affine-O-ample}.
\end{proof}

\begin{lemma}
\label{lemma-characterize-affine}
Let $X$ be a scheme. Suppose that there exist finitely many elements
$f_1, \ldots, f_n \in \Gamma(X, \mathcal{O}_X)$ such that
\begin{enumerate}
\item each $X_{f_i}$ is an affine open of $X$, and
\item the ideal generated by $f_1, \ldots, f_n$ in
$\Gamma(X, \mathcal{O}_X)$ is equal to the unit ideal.
\end{enumerate}
Then $X$ is affine.
\end{lemma}

\begin{proof}
Assume we have $f_1, \ldots, f_n$ as in the lemma.
We may write $1 = \sum g_i f_i$ for some $g_j \in \Gamma(X, \mathcal{O}_X)$
and hence it is clear that $X = \bigcup X_{f_i}$. (The $f_i$'s cannot
all vanish at a point.) Since each $X_{f_i}$
is quasi-compact (being affine) it follows that $X$ is quasi-compact.
Hence we see that $X$ is quasi-affine by
Lemma \ref{lemma-quasi-affine-O-ample} above.
Consider the open immersion
$$
j : X \to \Spec(\Gamma(X, \mathcal{O}_X)),
$$
see Lemma \ref{lemma-quasi-affine}. The inverse image of the standard open
$D(f_i)$ on the right hand side is equal to $X_{f_i}$ on the
left hand side and the morphism $j$ induces an isomorphism
$X_{f_i} \cong D(f_i)$, see
Lemma \ref{lemma-invert-f-affine}. Since the $f_i$ generate the unit ideal
we see that $\Spec(\Gamma(X, \mathcal{O}_X))
= \bigcup_{i = 1, \ldots, n} D(f_i)$. Thus $j$ is an isomorphism.
\end{proof}







